% ============================================================================
%  DisloCluster — 0-D solver and 0-D <-> 3-D coupling re-architecture
%
%  Build:  pdflatex solver_and_coupling_rearchitecture.tex   (twice, for refs)
%  Self-contained: no bibliography, no external figures.
% ============================================================================
\documentclass[11pt,a4paper]{article}

\usepackage[T1]{fontenc}
\usepackage[utf8]{inputenc}
\usepackage{lmodern}
\usepackage[margin=2.4cm]{geometry}
\usepackage{amsmath,amssymb}
\usepackage{booktabs}
\usepackage{array}
\usepackage{longtable}
\usepackage{graphicx}
\usepackage{xcolor}
\usepackage[colorlinks=true,linkcolor=blue!50!black,urlcolor=blue!50!black,
            citecolor=blue!50!black]{hyperref}

\newcommand{\code}[1]{\texttt{\small #1}}
\newcommand{\dpa}{\ensuremath{\mathrm{dpa}}}
\newcommand{\aloop}{\ensuremath{\langle a\rangle}}
\newcommand{\cloop}{\ensuremath{\langle c\rangle}}

\title{\bfseries Re-architecture of the DisloCluster 0-D Solver and of the
       0-D\,$\leftrightarrow$\,3-D Coupling}
\author{DisloCluster development note}
\date{8 August 2026}

\begin{document}
\maketitle

\begin{abstract}
\noindent
Three changes were made to the DisloCluster 0-D side and to its coupling with
the spatially resolved code. (i) The right-hand side was moved into a single
scalar-type-templated core so that the residual and an \emph{exact} Jacobian,
obtained by forward-mode automatic differentiation, are two instantiations of
one definition; SUNDIALS' difference-quotient Jacobian is thereby retired.
(ii) The six point-defect conservation accumulators were recognized as pure
quadratures and moved out of the implicit system into CVODES quadrature
variables, shrinking the Newton block from $19\times19$ to $13\times13$
standalone and to $9\times9$ under the operator split. (iii) The 0-D\,$\to$\,3-D
handoff, which previously collapsed the marched state to four scalars, was
replaced by a genuine per-node field exchange through MoDELib3's own
cluster-dynamics configuration block.

Measurement then redirected the effort: at $n=19$ the dense linear algebra was
never the bottleneck. Profiling with the integration window collapsed to zero
showed that per-case construction and destruction of SUNDIALS objects consumed
$1.31\,$s of thread time against $0.2\,$s of actual solving in a 256-point
coupling substep. Reusing one solver workspace per OpenMP thread cut the
end-to-end march wall time by a factor of $1.67$ at $q=1024$ quadrature points
--- more than the two originally requested changes combined. The report states
all three outcomes as measured, including the two that did not pay off in wall
time.

Two further sections document the coupling in operation.
Section~\ref{sec:fieldcost} measures the full 1--26 \dpa\ march on the
1~\textmu m single crystal in three configurations, separating the genuine
solver speedup ($1.95\times$ at identical work) from the cost of resolving a
field on 24\,115 nodes instead of averaging 25 line points to four scalars.
Section~\ref{sec:integrators} characterizes MoDELib3's own nodal immobile
integrator against the CVODE march --- correcting, along the way, the common
description of it as explicit Euler: it is a first-order semi-implicit scheme,
and it has to be, because the measured stiffness puts its step 46\,000 times
beyond the forward-Euler limit. Section~\ref{sec:provenance} then traces where
each scheme came from --- the immobile field does not exist at all in the
upstream MoDELib2-NNL --- and draws out the maintenance consequence of the two
different routes.
\end{abstract}

\tableofcontents

% ============================================================================
\section{Starting point}
% ============================================================================

The 0-D model integrates $N_{\mathrm{eq}}=19$ ordinary differential equations:
12 physical species, 6 point-defect conservation accumulators, and the evolving
network dislocation density $\rho_N$. In the operator-split coupling the four
mobile species $\mathbf C_M=[C_v,C_i,C_{2i},C_{3i}]$ are frozen at the value
supplied by the 3-D fast solve (\code{freeze\_mobile=1}), and the remaining
state is marched independently at every quadrature point, one case per point in
an OpenMP batch subprocess.

Before this work the solver had the following properties.

\begin{itemize}
\item \textbf{No Jacobian was supplied.} \code{CVodeSetJacFn} was never called,
      so SUNDIALS formed the $19\times19$ dense Jacobian by difference
      quotients, costing 19 extra residual evaluations per Jacobian setup.
\item \textbf{All 19 equations were implicit.} The accumulators were carried as
      ODE state, entered the Newton system and the error test, and were
      factorized in every LU.
\item \textbf{Every case built its own solver.} \code{integrate\_one} created
      and destroyed a \code{SUNContext}, state vector, dense matrix, linear
      solver and CVODE memory per invocation --- i.e.\ once per quadrature
      point per substep.
\item \textbf{The 0-D\,$\to$\,3-D carrier was four scalars.} The driver
      averaged the marched state over all points and wrote one number density
      and one mean radius per loop family into MoDELib's microstructure
      generator inputs.
\end{itemize}

% ============================================================================
\section{Structural observation: the accumulators are pure quadratures}
% ============================================================================

Rows 12--17 of the right-hand side are written but columns 12--17 are never
read: no equation in the system depends on an accumulator value. Inspection of
\code{rate\_equations.cpp} confirms that \code{y[12..17]} appears nowhere on the
right of any assignment. The Jacobian therefore has the block structure
\begin{equation}
J=\begin{pmatrix}
J_{ss} & 0 \\[2pt]
J_{as} & 0
\end{pmatrix},
\qquad
s=\{\text{species},\rho_N\},\quad a=\{\text{accumulators}\},
\end{equation}
so the accumulators are quadratures of the state trajectory,
$\mathbf a(t)=\int \mathbf g(\mathbf y_s(\tau))\,\mathrm d\tau$, and belong in
CVODES' quadrature module rather than in the Newton system.

Combining this with the operator split, which freezes the four mobile species,
gives the implicit block dimensions of Table~\ref{tab:dims}.

\begin{table}[htb]
\centering
\caption{Dimension of the implicit block actually factorized.}
\label{tab:dims}
\begin{tabular}{lccc}
\toprule
Configuration & Newton block $n$ & Dense LU $\propto n^3$ & Relative LU cost \\
\midrule
Legacy (all 19 implicit)                & 19 & 6859 & 1.00 \\
Reduced, standalone                     & 13 & 2197 & 0.32 \\
Reduced, \code{freeze\_mobile=1}        &  9 &  729 & 0.11 \\
\bottomrule
\end{tabular}
\end{table}

% ============================================================================
\section{Implementation}
% ============================================================================

\subsection{One templated core, two instantiations}

The physics was moved verbatim from \code{rate\_equations.cpp} into
\code{rate\_equations\_core.h} as
\begin{center}
\code{template <class T> void rhs\_core(const T* y, T* ydot, const Parameters\& P)}
\end{center}
instantiated with \code{T = double} for the residual and with
\code{T = Dual<N>} (\code{dual.h}) for the Jacobian. Forward-mode automatic
differentiation evaluates value and all $N$ directional derivatives in a
\emph{single} sweep: \code{sqrt}, \code{exp} and \code{cbrt} are evaluated once
and their derivatives propagated by a multiply, rather than the $N$ complete
residual evaluations a difference quotient needs.

Two properties motivated AD over a hand-written analytic Jacobian. First,
exactness: there is no perturbation size to choose and no truncation error.
Second, and more important for maintenance, the Jacobian cannot drift out of
sync with the residual --- hand-differentiating the coalescence Avrami gates and
the smoothstep minimum-size gate would have created a second copy of the physics
to keep aligned, on a code base that already maintains a Python mirror.

\paragraph{Bit-identity.} Floating-point addition is not associative, so the
transcription preserved the exact grouping of every expression. Where the
original wrote \code{k1 * k2 * C * D}, meaning $((k_1k_2)C)D$, the templated
core hoists $k_1k_2$ into a named \code{double} and keeps the remaining order.
This was verified rather than asserted: a reference binary was built from the
git-HEAD sources and the outputs compared byte for byte
(Section~\ref{sec:results}).

\subsection{Reduced implicit block}

With \code{--reduced=1} the state vector holds only the implicit components and
the accumulators are advanced by \code{CVodeQuadInit}. Quadrature error control
is left at its default (off), so the accumulators ride the step sequence the
state equations demand instead of constraining it. The 19-column output
contract is unchanged: each row is reassembled from the reduced state, the
frozen mobile values and the quadrature vector.

\subsection{Reusable per-thread solver workspace}
\label{sec:workspace}

This change was not planned; it was the result of profiling. Running a 256-case
frozen-mobile batch with the integration window collapsed to
$t_{\mathrm{end}}=t_{\mathrm{begin}}(1+10^{-12})$ --- so that only process
spawn, case-file parsing and output formatting remain --- gave the decomposition
in Table~\ref{tab:overhead}.

\begin{table}[htb]
\centering
\caption{Cost decomposition of one 256-case coupling substep, before the
workspace change. ``cpu-solve'' is the sum of per-case integration wall times
over all threads.}
\label{tab:overhead}
\begin{tabular}{lrr}
\toprule
Configuration & wall [s] & cpu-solve [s] \\
\midrule
Null integration (setup/teardown only) & 0.172 & 1.310 \\
1 \dpa, $19\times19$ difference-quotient & 0.136 & 1.608 \\
1 \dpa, $19\times19$ AD                  & 0.129 & 1.491 \\
\bottomrule
\end{tabular}
\end{table}

Constructing and destroying the SUNDIALS objects cost $1.31\,$s of thread time
against roughly $0.2\,$s of genuine integration: the per-case setup was
consuming about six times what the ODE solve did. A \code{Workspace} is now
created once per OpenMP thread and reused across every case that thread handles
via \code{CVodeReInit} / \code{CVodeQuadReInit}, rebuilt only when a case needs
a different shape. In a 1024-point march this reduces 1024 solver
constructions per substep to 24, one per thread.

\subsection{Interface}

Three flags were added; both numerical flags default to \emph{off} so that a
standalone 0-D run remains bit-identical to the previous solver.
\code{modelib\_coupling.build\_immobile\_cases} turns them on for the march,
where they are measurably better.

\begin{center}
\begin{tabular}{lp{10.2cm}}
\toprule
Flag & Effect \\
\midrule
\code{--reduced=1}      & Accumulators as CVODES quadrature; Newton block 13
                          (or 9 with \code{freeze\_mobile}). CVODE only. \\
\code{--analytic\_jac=1}& Exact dense Jacobian by forward AD. Dense linear
                          solver only. \\
\code{--stats=1}        & Emit a \code{\# STATS} comment line with the
                          integrator counters. Non-numeric, so the existing
                          Python row parsers skip it. \\
\bottomrule
\end{tabular}
\end{center}

% ============================================================================
\section{Verification}
% ============================================================================

\subsection{The AD Jacobian}

A standalone tool (\code{jac\_check.cpp}) compares the AD Jacobian against a
central-difference Jacobian of the \emph{same} core, differencing column $j$
with a relative step $h_j=\varepsilon_{\mathrm{rel}}|y_j|$ because the state
spans some 34 decades. It was run at 20 states sampled along a converged
trajectory from $10^{-8}$ to $100$ \dpa.

Two classes of entry are excluded, for reasons belonging to the finite
difference and not to the AD:

\begin{itemize}
\item \textbf{FD-unresolvable entries.} A central difference can only see an
      entry whose contribution $J_{ij}h_j$ rises above the floating-point
      granularity of $f_i$. At the initial condition, $\mathrm dC_i/\mathrm dt$
      is pinned near $G_i\approx10^{-7}$ while its $C_v$ derivative moves it by
      ${\sim}10^{-28}$, i.e.\ $10^{-5}$ of one unit in the last place. The
      difference quotient returns exactly zero there at \emph{any} step size;
      AD returns the exact value. Entries are compared only where the signal
      exceeds 100 ULP. Note that this is precisely the blindness SUNDIALS'
      difference-quotient Jacobian suffers, so the AD Jacobian is strictly more
      informative in this region.
\item \textbf{Components on the concentration floor}, where \code{fl()} makes
      the residual non-smooth: AD returns the one-sided derivative while a
      central difference straddles the kink.
\end{itemize}

The result is in Table~\ref{tab:jac}. The discrepancy falls by two decades per
decade of $\varepsilon_{\mathrm{rel}}$ across the usable range, the $O(h^2)$
signature of a correct derivative with the finite difference's own truncation
error dominating.

\begin{table}[htb]
\centering
\caption{AD versus central-difference Jacobian, worst row-relative discrepancy
over all compared entries. Representative states from the 20 tested.}
\label{tab:jac}
\begin{tabular}{rrrrr}
\toprule
Dose [\dpa] & $\varepsilon_{\mathrm{rel}}=10^{-5}$ & $10^{-6}$ & $10^{-7}$
            & entries compared \\
\midrule
$1.0\times10^{-8}$ & $6.50\times10^{-12}$ & $8.84\times10^{-11}$ & $1.24\times10^{-9}$ & 6 \\
$1.3\times10^{-6}$ & $8.49\times10^{-8}$  & $7.96\times10^{-7}$  & $2.05\times10^{-6}$ & 108 \\
$1.6\times10^{-4}$ & $8.29\times10^{-8}$  & $9.97\times10^{-7}$  & $1.83\times10^{-5}$ & 117 \\
$2.1\times10^{-2}$ & $4.40\times10^{-7}$  & $4.71\times10^{-6}$  & $5.16\times10^{-5}$ & 133 \\
$7.8\times10^{-1}$ & $1.55\times10^{-7}$  & $9.30\times10^{-7}$  & $9.92\times10^{-6}$ & 133 \\
$8.9\times10^{0}$  & $3.09\times10^{-9}$  & $4.28\times10^{-8}$  & $1.30\times10^{-6}$ & 107 \\
$1.0\times10^{2}$  & $8.04\times10^{-9}$  & $5.17\times10^{-8}$  & $5.71\times10^{-7}$ & 107 \\
\midrule
\multicolumn{5}{l}{Worst over all 20 states (best $\varepsilon$ per state):
 $4.40\times10^{-7}$; states failing at every $\varepsilon$: 0/20.} \\
\bottomrule
\end{tabular}
\end{table}

\subsection{Bit-identity of the untouched path}

The legacy configuration of the rebuilt solver reproduces the git-HEAD
reference binary \emph{byte for byte} --- 1000 output rows $\times$ 19 columns
in the standalone run, and the endpoint state of all 256 points $\times$ 5
substeps in the coupling march. The refactor into a templated core therefore
changed no arithmetic.

% ============================================================================
\section{Results}
\label{sec:results}
% ============================================================================

All runs use \code{input/Zr\_input\_parameters.xlsx} at $T=573$~K,
$G=10^{-7}$~\dpa/s, on 24 cores. Because wall-clock times on a shared machine
are noisy, exact operation counts are reported alongside and are the primary
evidence:
\begin{align*}
n_{\mathrm{core}} &= \text{residual-core evaluations} =
  \code{nfe} + \code{nfeLS} + \code{nqmiss},\\
n_{\mathrm{LU}}   &= \text{dense factorizations (\code{nsetups})}, \qquad
\text{LU flops} \propto n_{\mathrm{LU}}\,n^{3}/3 .
\end{align*}

\subsection{Standalone 0-D run}

$t=10^{-1}\ldots10^{9}$~s (100 \dpa), 1000 output points, rtol $=10^{-6}$,
atol $=10^{-20}$. Accuracy is measured against a tight reference
(rtol $=10^{-10}$, atol $=10^{-24}$), compared per column only where the
reference exceeds $10^{-8}$ of that column's own peak --- a plain elementwise
relative error is meaningless on a state spanning 34 decades whose early
transient sits below atol.

\begin{table}[htb]
\centering
\caption{Standalone 0-D run. ``bit=ref'' is byte-for-byte agreement with the
git-HEAD binary.}
\label{tab:phaseA}
\begin{tabular}{lrrrrrrc}
\toprule
Configuration & wall & solve & \code{nst} & $n_{\mathrm{core}}$
              & $n_{\mathrm{LU}}$ & LU flops & bit=ref \\
              & [ms] & [ms] & & & & & \\
\midrule
git HEAD $19\times19$ DQ  & 70.7 & --   & --   & --   & --  & --                  & --  \\
Baseline $19\times19$ DQ  & 69.3 & 46.7 & 2797 & 6129 & 466 & $1.07\times10^{6}$  & yes \\
AD Jac.\ $19\times19$ AD  & 65.5 & 46.9 & 2483 & 4171 & 429 & $9.81\times10^{5}$  & no  \\
Reduced $13\times13$ DQ   & 64.5 & 42.3 & 2564 & 7823 & 423 & $3.10\times10^{5}$  & no  \\
Reduced\,+\,AD $13\times13$ & 63.4 & 43.7 & 2811 & 7415 & 458 & $3.35\times10^{5}$ & no \\
\bottomrule
\end{tabular}
\end{table}

\begin{table}[htb]
\centering
\caption{Standalone accuracy against the tight reference (max relative error).}
\label{tab:phaseAerr}
\begin{tabular}{lrrr}
\toprule
Configuration & species & accumulators & $\rho_N$ \\
\midrule
git HEAD / baseline $19\times19$ DQ & $1.41\times10^{-4}$ & $2.15\times10^{-5}$ & $1.60\times10^{-5}$ \\
AD Jac.\ $19\times19$               & $2.18\times10^{-4}$ & $4.90\times10^{-5}$ & $3.16\times10^{-5}$ \\
Reduced $13\times13$ DQ             & $2.44\times10^{-4}$ & $5.25\times10^{-5}$ & $3.30\times10^{-5}$ \\
Reduced\,+\,AD $13\times13$         & $8.79\times10^{-5}$ & $2.18\times10^{-5}$ & $1.30\times10^{-5}$ \\
\bottomrule
\end{tabular}
\end{table}

The AD Jacobian removes 32\% of the residual evaluations
($6129\to4171$) and 11\% of the steps. The reduction cuts LU work by 71\% but
\emph{adds} residual evaluations, for the reason given in
Section~\ref{sec:memo}. All four configurations sit within a factor of three of
each other in accuracy, at the level expected from different step sequences at
the same tolerance; none is meaningfully more accurate than another.

\subsection{Operator-split coupling march}

A 5 \dpa\ dose step in 5 substeps, seeded at 1 \dpa, with the mobile field
depleted toward a sink across the point set. Table~\ref{tab:phaseB} gives
$q=1024$; the $q=64$ and $q=256$ rows scale proportionally.

\begin{table}[htb]
\centering
\caption{Coupling march, $q=1024$ quadrature points, 5 substeps.
``warn'' counts CVODE roundoff warnings (see Section~\ref{sec:pathology});
error is against the git-HEAD result.}
\label{tab:phaseB}
\begin{tabular}{lrrrrrrr}
\toprule
Configuration & wall [s] & cpu-solve [s] & \code{nst} & $n_{\mathrm{core}}$
              & $n_{\mathrm{LU}}$ & LU flops & warn \\
\midrule
git HEAD $19\times19$ DQ    & 2.379 & --     & --     & --     & --    & --                 & 2048 \\
Reuse $19\times19$ DQ       & 1.424 & 12.03  & 175771 & 366170 & 69860 & $1.60\times10^{8}$ & 2048 \\
Reuse\,+\,AD $19\times19$   & 1.435 & 12.98  & 175890 & 243682 & 69921 & $1.60\times10^{8}$ & 2048 \\
Reduced $9\times9$ DQ       & 1.568 & 13.99  & 143651 & 400404 & 54959 & $1.34\times10^{7}$ & 0 \\
Reduced\,+\,AD $9\times9$   & 1.572 & 13.68  & 143654 & 349626 & 54957 & $1.34\times10^{7}$ & 0 \\
\bottomrule
\end{tabular}
\end{table}

\begin{table}[htb]
\centering
\caption{March wall time versus point count. Speedup is against the git-HEAD
binary at the same $q$.}
\label{tab:scaling}
\begin{tabular}{lrrrrrr}
\toprule
& \multicolumn{2}{c}{$q=64$} & \multicolumn{2}{c}{$q=256$}
& \multicolumn{2}{c}{$q=1024$} \\
\cmidrule(lr){2-3}\cmidrule(lr){4-5}\cmidrule(lr){6-7}
Configuration & [s] & speedup & [s] & speedup & [s] & speedup \\
\midrule
git HEAD                  & 0.254 & 1.00 & 0.692 & 1.00 & 2.379 & 1.00 \\
Reuse $19\times19$ DQ     & 0.214 & 1.19 & 0.482 & 1.44 & 1.424 & \textbf{1.67} \\
Reuse\,+\,AD $19\times19$ & 0.211 & 1.20 & 0.468 & 1.48 & 1.435 & 1.66 \\
Reduced\,+\,AD $9\times9$ & 0.217 & 1.17 & 0.502 & 1.38 & 1.572 & 1.51 \\
\bottomrule
\end{tabular}
\end{table}

Agreement with the git-HEAD march is $1.3\times10^{-5}$ or better in every
species at every $q$ --- the expected splitting-consistent difference between
step sequences, not a change of answer.

\subsection{Where the improvement actually came from}

Stated plainly, and against the expectation that motivated the work:

\begin{itemize}
\item \textbf{Workspace reuse is the whole wall-clock gain:} $1.67\times$ at
      $q=1024$, and it grows with $q$ because the per-case setup it removes
      scales with the point count while the per-thread construction does not.
\item \textbf{The AD Jacobian delivers its target metric:} residual evaluations
      fall by 33\% ($366\,170\to243\,682$) and difference-quotient evaluations
      go to zero. Wall time is neutral, because one \code{Dual<19>} sweep costs
      about what the 19 scalar evaluations it replaces cost. Its durable value
      is exactness, the elimination of the FD-blind entries of
      Table~\ref{tab:jac}, and a Jacobian that cannot fall out of step with the
      residual.
\item \textbf{The $9\times9$ reduction delivers its target metric --- a 92\%
      cut in LU flops --- but LU was never the bottleneck at $n=19$.} It also
      reduces the step count by 18\%. It nonetheless loses slightly in wall
      time, because of the quadrature cost analyzed next.
\end{itemize}

% ============================================================================
\subsection{A step-size pathology the reduction cures}
\label{sec:pathology}
% ============================================================================

The \code{warn} column of Table~\ref{tab:phaseB} is not cosmetic. The legacy
path emits, on every substep of every point,
\begin{center}
\code{[WARNING] Internal t = 9.88525e+06 and h = 2.10129e-12 are such that
      t + h = t on the next step}
\end{center}
2048 of them in a 1024-point, 5-substep march.

The cause is the accumulators. Each substep restarts at $t\approx10^{7}$~s with
the accumulators reset to zero, and with atol $=10^{-20}$ their error weights
$w_i=1/(\mathrm{rtol}\,|y_i|+\mathrm{atol})$ are ${\sim}10^{20}$. CVODE's
initial-step heuristic uses the weighted norm of $\dot{\mathbf y}$; with
$\dot a\sim10^{-7}$ this gives $\|\dot{\mathbf y}\|_{\mathrm{WRMS}}\sim10^{13}$
and hence $h_0\sim10^{-13}$, while $\varepsilon t\approx2\times10^{-10}$. The
first step is below the representable resolution of $t$ and is wasted.

Removing the accumulators from the error test eliminates this outright: zero
warnings, and 18\% fewer internal steps. This is the reduction's most defensible
benefit --- it is a robustness property, independent of problem size, whereas
the LU saving is not.

A second, subtler effect works the other way. The error test is a root-mean
square over components, $\|e\|_{\mathrm{WRMS}}=\sqrt{N^{-1}\sum_i(e_iw_i)^2}$,
so removing six well-behaved components makes the norm $\sqrt{19/9}=1.45$ times
more sensitive per remaining component. The 19-equation system was being
flattered by six dummy components diluting its error norm; the reduced system
is held to a genuinely tighter standard at the same nominal tolerance.

% ============================================================================
\subsection{The quadrature memo, and why the reduction does not win on time}
\label{sec:memo}
% ============================================================================

CVODES calls the quadrature right-hand side once per completed step. Because
that call needs the same core evaluation the residual has just performed, the
implementation memoizes the accumulator rates on the exact $(t,\mathbf y)$ bit
pattern and reuses them on a match.

The measured hit rate is poor: 1 hit against 2564 misses in the standalone run.
CVODE's corrector evaluates the residual at its Newton iterates, and the
accepted step point $\mathbf y_n$ is generally not one of them, so the memo
misses and a second full core sweep is paid per step. That is exactly the
${+}28\%$ in $n_{\mathrm{core}}$ visible in Table~\ref{tab:phaseA}, and it
cancels the LU saving.

The identified remedy, not implemented here, is to keep the accumulators in the
state vector but neutralize their error weights with a vector \code{atol}
(\code{CVodeSVtolerances} with atol $=10^{300}$ on those six components). That
would retain the step-size cure of Section~\ref{sec:pathology} at \emph{zero}
extra core evaluations, at the cost of a $15\times15$ rather than $9\times9$
block --- which, on this evidence, is the better trade. It is recorded as the
next step rather than claimed as a result.

% ============================================================================
\section{The 0-D $\leftrightarrow$ 3-D field handoff}
% ============================================================================

\subsection{What was wrong}

The previous 0-D\,$\to$\,3-D carrier
(\code{modelib\_fem.write\_sink\_field}) wrote four scalars per dose step --- one
number density and one mean radius per loop family --- into
\code{inputFiles/aLoopsDensity.txt} and \code{inputFiles/frankLoopsDensity.txt}.
Those files drive MoDELib's \emph{microstructure generator}, which places
discrete loops at a uniform target density. The driver therefore had to collapse
the marched state to its volume average before handing it over:
\begin{center}
\code{A = np.mean(np.asarray(Y), axis=0)}
\end{center}
Every bit of spatial structure the slow march produced was discarded on the way
out. In the reverse direction the mobile field was read from
\code{F/quadrature\_*.txt} through a column map whose source comment conceded it
had never been checked against a real run.

\subsection{What MoDELib3 actually offers}

MoDELib3 does not need a scalar. \code{ClusterDynamicsFEM} already carries the
immobile population as a finite-element \emph{field} --- \code{immobileClusters},
with $\code{iSize}=8$ components per node --- and \code{ImmobileSinks} builds the
sink strength from the local value at every quadrature point.
\code{ClusterDynamics::output} writes that field, together with the
$\code{mSize}=4$ mobile components, into the CD block of
\code{evl/evl\_<N>.txt}. Decisively, \code{ClusterDynamicsFEM::\allowbreak
initializeConfiguration} \emph{reads it back}:

\begin{quote}\small\ttfamily
if(size\_t(configIO.cdMatrix().size())==mobileClusters.gSize()+immobileClusters.gSize())\\
\{\\
\hspace*{1em}mobileClusters   = configIO.cdMatrix().block(0,0,nNodes,mSize)\ldots\\
\hspace*{1em}immobileClusters = configIO.cdMatrix().block(0,mSize,nNodes,iSize)\ldots\\
\}
\end{quote}

The CD block is therefore a read/write field channel in both directions, on the
node set the two codes already share. \code{evl/cdNodes.txt} gives the node
coordinates in the same row order, so no search or interpolation is required.

\subsection{The mapping}

\begin{table}[htb]
\centering
\caption{CD block layout, \code{ClusterDynamicsParameters}: mSize $=4$,
iSize $=8$. \code{immobileSpeciesVector = -1 1 1 1} makes family 0 the basal
\cloop\ vacancy family and families 1--3 the three prism \aloop\ variants.}
\begin{tabular}{clll}
\toprule
Columns & Block & Components & Units \\
\midrule
0--3   & mobile   & $C_v,\ C_i,\ C_{2i},\ C_{3i}$ & atom fraction \\
4--7   & immobile & $n_c,\ n_{a1},\ n_{a2},\ n_{a3}$ & per $b^3$ \\
8--11  & immobile & $c_c,\ c_{a1},\ c_{a2},\ c_{a3}$ & atom fraction \\
\bottomrule
\end{tabular}
\end{table}

MoDELib works in Burgers-vector units, so its nodal number density is per $b^3$
while the 0-D loop density is an atom fraction. \code{ImmobileSinks} floors them
as $n_{\mathrm{floor}}=C_{\mathrm{floor}}/\omega$ and
$c_{\mathrm{floor}}=C_{\mathrm{floor}}$, which fixes the conversion:
\begin{equation}
n_{\mathrm{MoDELib}}=\frac{C_{\mathrm{0D}}}{\omega},\qquad
c_{\mathrm{MoDELib}}=c_{\mathrm{0D}},\qquad
\omega=\frac{\code{atomicVolume\_SI}}{\code{b\_SI}^{3}}=0.355111\ b^{3}.
\end{equation}
The 0-D aligned/non-aligned pair is collapsed and the \aloop\ total distributed
over the three prism variants by the resolved-stress weights --- the same
Option-A reduction \code{modelib\_export} performs, but applied \emph{pointwise}
rather than to a volume average. The mobile direction is a pure slice: MoDELib's
four mobile species \emph{are} the 0-D mobile species, in the same order and
units.

\subsection{Verification on the 3-D verification case}

Tested against \code{MoDELib3/tutorials/zrmicro\_coupled}, the 1~\textmu m
single-crystal case with 24\,115 CD nodes, at \code{evl\_30.txt}.

\begin{table}[htb]
\centering
\caption{Field-handoff verification, 24\,115 CD nodes.}
\label{tab:field}
\begin{tabular}{lp{9.3cm}}
\toprule
Check & Result \\
\midrule
CD block rewrite round-trip & identical, $24\,115\times12$; header and all
                              dislocation records preserved \\
0-D\,$\to$\,MoDELib\,$\to$\,0-D immobile & max relative error
                              $4.47\times10^{-16}$ (machine precision) \\
Mobile field read back      & $C_v$ spans $8.46\times10^{-15}$ to
                              $6.70\times10^{-6}$, a dynamic range of
                              $7.9\times10^{8}$ \\
Spatial structure           & $\mathrm{corr}(\log_{10}C_v,\ \text{distance to
                              nearest face}) = +0.63$ \\
March on the real node set  & 24\,115 nodes over 1 \dpa\ in 9.5~s, 0 failures \\
Field written back          & $n_{a1}$ spans a factor 4.4 across the domain \\
\bottomrule
\end{tabular}
\end{table}

The information content of the handoff rises from 4 numbers to
$24\,115\times12=289\,380$ numbers, and the boundary-layer structure that the
Dirichlet condition imposes on the mobile field --- eight decades in $C_v$,
correlating with distance to the surface --- now reaches the slow march instead
of being averaged away.

\paragraph{A precision detail.} Restoring the aligned/non-aligned split on the
return trip initially lost eight significant digits on the minor partner. The
fractions were formed as $f$ and $1-f$; when the pair is lopsided (and
$C_{iL}\gg C_{aiL}$ is normal at low dose) $f$ rounds to $1-\epsilon$ and $1-f$
retains only the bits surviving the cancellation. Forming both fractions
directly from the pair restores full precision, giving the
$4.47\times10^{-16}$ of Table~\ref{tab:field}.

\subsection{Notebook wiring}

\code{coupled\_0d\_3d\_ZrMicro.ipynb} gains a \code{"field"} fast-solve backend,
selected automatically when a CD field is present. In that mode the march runs
on the CD node set rather than the synthetic 1-D line, seeds the immobile state
from the 3-D field itself, and after each dose interval writes
\code{evl/evl\_coupled\_<dose>dpa.txt} --- a complete configuration, with the
dislocation records copied through unchanged, that MoDELib can be restarted
from. The previous scalar path and the analytic placeholder both remain
available and unmodified.

Because MoDELib reads its configuration at start-up, the handback takes effect
on a \emph{restart} from the written file. That is the rhythm the operator split
requires anyway (fast solve $\to$ slow march $\to$ fast solve), but it does mean
the coupling is restart-driven rather than in-process: at the default six
snapshots that is six DDomp launches, which is unremarkable; at hundreds of dose
intervals it would argue for driving MoDELib through \code{pyMoDELib} instead.

% ---------------------------------------------------------------------------
\subsection{Measured cost of the field upgrade}
\label{sec:fieldcost}
% ---------------------------------------------------------------------------

The comparison needs care, because the old path never ran on the mesh at all.
\code{RUN["n\_points"] = 25}: it marched a synthetic 1-D line of 25 points bearing
no relation to the geometry, while the field path marches the 24\,115 CD nodes
--- 965 times as many. A same-work control is therefore included: the git-HEAD
solver driven over the \emph{same} 24\,115 nodes, which is what the old solver
would have cost to do the new job.

\paragraph{Test.} The notebook's default schedule --- 1 to 26 \dpa, five
intervals, \code{substeps\_per\_interval = 20}, i.e.\ 100 batch subprocesses ---
run to completion in all three configurations, seeded at 1 \dpa\ from the 0-D
reference and, for the field configurations, from the CD field itself.

\begin{table}[htb]
\centering
\caption{Full 1--26 \dpa\ coupling march, 100 batch subprocesses, 24 cores.}
\label{tab:march1um}
\begin{tabular}{llrrrr}
\toprule
& Configuration & $q$ & wall [s] & s/substep & ms/point/substep \\
\midrule
A & Old as configured (git HEAD, 1-D line) & 25      & 4.3    & 0.04  & 1.715 \\
B & Old solver at the new workload         & 24\,115 & 1117.4 & 11.17 & 0.463 \\
C & New: $9\times9$ + AD + reuse           & 24\,115 & 574.2  & 5.74  & 0.238 \\
\bottomrule
\end{tabular}
\end{table}

Two distinct readings, and both are needed to avoid misrepresenting the change:

\begin{itemize}
\item \textbf{Same work, B versus C: $1.95\times$ faster.} That is the honest
      solver speedup --- identical points, identical schedule. It exceeds the
      $1.67\times$ measured at $q=1024$ in Table~\ref{tab:scaling} because the
      per-case solver construction the old binary pays scales with $q$ while the
      new binary's per-thread construction does not: 24\,115 solver builds per
      substep become 24.
\item \textbf{Different work, A versus C: $134\times$ more wall clock for
      $965\times$ the points.} The old configuration is faster in absolute terms
      only because it was doing $1/965$ of the work and emitting eight numbers.
      Per point per substep the new path is $7.2\times$ cheaper
      ($1.715\to0.238$~ms), part solver improvement and part fixed
      process-and-parse overhead amortizing over far more cases.
\end{itemize}

For scale: the 3-D solve itself (\code{clean\_run.sh} $\to$ DDomp) takes hours.
Against that, neither 4~s nor 10~min makes the march the pacing item of a
coupled run.

\paragraph{Consistency.} Over the full 1--26 \dpa\ march, configurations B and C
agree to $9.05\times10^{-7}$ in the species. That is \emph{tighter} than the
$8.2\times10^{-4}$ measured after the first 5 \dpa\ interval alone, and the
direction is informative: the early transient difference between two step
sequences washes out rather than accumulating, because the loop populations are
approaching a quasi-steady state and the dynamics are attracting. The largest
disagreements sit at the stiff near-boundary nodes, where the frozen mobile field
is eight decades below the interior value. It is a tolerance effect, not a change
of answer.

% ============================================================================
\section{Integrator comparison: MoDELib's nodal update versus the CVODE march}
\label{sec:integrators}
% ============================================================================

The two codes advance the \emph{same} immobile physics --- eight components per
node --- with entirely different schemes. This section characterizes both, since
the difference is a systematic source of 0-D/3-D discrepancy that is independent
of the physics non-correspondences already documented in \code{CLAUDE.md}.

\subsection{What MoDELib actually does}

It is worth stating precisely, because the scheme is often described as explicit
Euler and it is not. \code{ClusterDynamicsFEM::solveImmobileClusters} loops over
nodes and, inside each node, over \code{nSub = 20} substeps of
$\mathrm dt = \mathrm{dtMax}/20$, applying

\begin{quote}\small\ttfamily
// implicit for the linear losses, explicit for the sources\\
const double nNew((n(k)+dt*nucRate)/(1.0+dt*lossN));\\
const double cNew((c(k)+dt*(gdot(k)+Gk+clusShare(k)*clusC(pol)-contentSink))/(1.0+dt*lossC));
\end{quote}

This is a first-order \textbf{semi-implicit (IMEX)} update: the sources
(nucleation, absorbed flux, cascade seed) are explicit, while the linear loss
coefficients --- coalescence, $1/\tau_{vL}$, and the thermal-emission rate
recast as a first-order coefficient --- sit in the denominator. The source
comment records the reason: \emph{``$\tau_{vL}$ is short compared with a dose
step \ldots an explicit update would give $n(1-\mathrm dt/\tau)<0$ and collapse
the vacancy-loop density.''}

With $G=10^{-7}$~\dpa/s and \code{dtMax = 6.95868964e19} in MoDELib units (time
unit $1.4371\times10^{-13}$~s), one dose step is $10^{7}$~s and the immobile
substep is
\begin{equation}
\mathrm dt_{\mathrm{MoDELib}} = \frac{10^{7}\ \mathrm s}{20}
 = 5\times10^{5}\ \mathrm s .
\end{equation}

\subsection{Test method}

Three instruments, all exercising the \emph{same} \code{rate\_equations\_core.h}
right-hand side so that no model difference contaminates the comparison:

\begin{enumerate}
\item \textbf{Spectrum.} The frozen-mobile $9\times9$ immobile Jacobian is formed
      at states sampled along a converged trajectory and its eigenvalues taken.
      The forward-Euler stability limit is $\mathrm dt<2/|\mathrm{Re}\,\lambda|_{\max}$.
\item \textbf{Fixed-step explicit reference.} A \code{--euler\_nsub=N} mode was
      added to the solver: $N$ forward-Euler substeps per output interval on the
      same core, with the per-substep positivity clamp MoDELib applies. This is
      the \emph{control} --- what the update would do without the implicit loss
      treatment --- not a model of MoDELib's scheme.
\item \textbf{Reference solution.} CVODE at \code{rtol=1e-12}, \code{atol=1e-26}
      with the AD Jacobian, against which every other result is measured.
\end{enumerate}

Two nodes of the 1~\textmu m case bracket the range: the interior node of
maximum $C_v$ and the near-boundary node of minimum $C_v$.

\subsection{Stability}

\begin{table}[htb]
\centering
\caption{Spectrum of the frozen-mobile $9\times9$ immobile Jacobian.}
\label{tab:spectrum}
\begin{tabular}{rrrrr}
\toprule
Dose [\dpa] & $|\mathrm{Re}\,\lambda|_{\max}$ [s$^{-1}$]
& $\mathrm dt_{\mathrm{stab}}=2/|\mathrm{Re}\,\lambda|_{\max}$
& stiffness ratio & $\mathrm dt_{\mathrm{MoDELib}}/\mathrm dt_{\mathrm{stab}}$ \\
\midrule
 0.99 & $1.854\times10^{-1}$ & 10.8 s & $1.5\times10^{5}$ & 46\,338 \\
 2.96 & $1.853\times10^{-1}$ & 10.8 s & $5.7\times10^{7}$ & 46\,335 \\
 5.91 & $1.853\times10^{-1}$ & 10.8 s & $4.0\times10^{5}$ & 46\,335 \\
11.16 & $1.853\times10^{-1}$ & 10.8 s & $1.2\times10^{5}$ & 46\,335 \\
15.78 & $1.853\times10^{-1}$ & 10.8 s & $3.1\times10^{4}$ & 46\,335 \\
21.05 & $1.853\times10^{-1}$ & 10.8 s & $1.8\times10^{7}$ & 46\,335 \\
26.52 & $1.853\times10^{-1}$ & 10.8 s & $9.9\times10^{4}$ & 46\,335 \\
\bottomrule
\end{tabular}
\end{table}

The fast eigenvalue is essentially dose-independent at
$0.185$~s$^{-1}$ --- the vacancy-loop dissolution channel, $n_{vL,\mathrm{nuc}}/\tau_{vL}$
--- so a forward-Euler step must satisfy $\mathrm dt<10.8$~s. MoDELib runs at
$5\times10^{5}$~s, \textbf{46\,000 times the explicit limit}. Its implicit loss
treatment is therefore not a refinement; it is what makes the scheme usable at
all.

Table~\ref{tab:collapse} shows the consequence directly, from the fixed-step
explicit control at MoDELib's own substep count.

\begin{table}[htb]
\centering
\caption{One 1 \dpa\ step at the near-boundary node: a genuine forward Euler at
$\mathrm dt=5\times10^{5}$~s annihilates the entire \cloop\ vacancy-loop family,
exactly the failure the MoDELib source comment anticipates. The \aloop\ families,
which carry no fast decay channel, are unaffected.}
\label{tab:collapse}
\begin{tabular}{lrrrr}
\toprule
Species & initial & CVODE reference & explicit Euler, $N=20$ & rel.\ error \\
\midrule
$C_{iL}$   & $2.53\times10^{-3}$ & $5.69\times10^{-3}$ & $5.69\times10^{-3}$ & $1.8\times10^{-11}$ \\
$C_{aiL}$  & $1.60\times10^{-3}$ & $3.44\times10^{-3}$ & $3.44\times10^{-3}$ & $2.9\times10^{-11}$ \\
$C_{iL,i}$ & $2.16\times10^{-1}$ & $2.79\times10^{-1}$ & $2.79\times10^{-1}$ & $2.2\times10^{-10}$ \\
$C_{aiL,i}$& $1.37\times10^{-1}$ & $1.73\times10^{-1}$ & $1.73\times10^{-1}$ & $2.3\times10^{-10}$ \\
\midrule
$C_{vL}$   & $3.48\times10^{-10}$ & $3.48\times10^{-10}$ & $\mathbf{1.0\times10^{-20}}$ (floor) & $1.00$ \\
$C_{avL}$  & $2.02\times10^{-10}$ & $2.02\times10^{-10}$ & $\mathbf{1.0\times10^{-20}}$ (floor) & $1.00$ \\
$C_{vL,v}$ & $7.21\times10^{-4}$  & $7.21\times10^{-4}$  & $\mathbf{1.0\times10^{-20}}$ (floor) & $1.00$ \\
$C_{avL,v}$& $4.20\times10^{-4}$  & $4.20\times10^{-4}$  & $\mathbf{1.0\times10^{-20}}$ (floor) & $1.00$ \\
\bottomrule
\end{tabular}
\end{table}

Both schemes actually in use are unconditionally stable on the stiff channel,
but by different mechanisms. MoDELib's amplification factor
$1/(1+\mathrm dt\,\lambda)\in(0,1)$ for every $\mathrm dt>0$ and $\lambda>0$:
L-stable on the linear part, and positivity-preserving by construction. CVODE
BDF is A-stable at orders 1--2 and stiffly stable to order 5, and in addition
selects its own step from a local error estimate.

\subsection{Accuracy and speed}

\begin{table}[htb]
\centering
\caption{One 1 \dpa\ dose step, interior node ($C_v=6.70\times10^{-6}$), against
the \code{rtol=1e-12} reference. ``core evals'' counts evaluations of the shared
right-hand side.}
\label{tab:eulerbdf}
\begin{tabular}{lrrrrr}
\toprule
Integrator & substeps & $\mathrm dt$ [s] & core evals & max rel.\ error & $t$ [ms] \\
\midrule
Explicit Euler (MoDELib's \code{nSub}) & 20    & $5.0\times10^{5}$ & 20    & $9.27\times10^{4}$ & 0.14 \\
Explicit Euler                          & 200   & $5.0\times10^{4}$ & 200   & $9.27\times10^{3}$ & 0.27 \\
Explicit Euler                          & 2000  & $5.0\times10^{3}$ & 2000  & $9.26\times10^{2}$ & 1.18 \\
Explicit Euler                          & 20000 & $5.0\times10^{2}$ & 20000 & $9.17\times10^{1}$ & 13.71 \\
\midrule
CVODE BDF $9\times9$ + AD, \code{rtol=1e-6} & 161 (adaptive) & $6.2\times10^{4}$ & 366 & $\mathbf{2.10\times10^{-6}}$ & 1.50 \\
\bottomrule
\end{tabular}
\end{table}

At the near-boundary node CVODE needs only 37 steps and 84 core evaluations for
$1.27\times10^{-8}$, the frozen mobile field there being effectively zero.

The explicit control converges cleanly at first order --- a factor 10 in error
per factor 10 in step count --- which confirms the order; the enormous constant
reflects that it is simultaneously unstable.

\paragraph{What can and cannot be claimed for MoDELib's scheme.} Its IMEX
splitting was not reimplemented on the ZrMicro equations, so no error
\emph{number} is quoted for it here. What is structural, and visible in the code
itself, is that it is \textbf{first-order accurate with no error control}: no
local error estimate, no step rejection, no adaptivity, and 20 fixed substeps per
dose step whatever the local stiffness. Against that, CVODE selects its own order
(up to 5) and step from a local error estimate at \code{rtol}$=10^{-6}$.

\paragraph{Speed.} Per node per dose step, MoDELib performs 20 nodal updates of
eight components with no Jacobian, no linear solve and no error estimate; the
comparable explicit control costs $0.14$~ms against CVODE's $1.50$~ms interior
and $0.76$~ms near the boundary. So MoDELib's immobile update is roughly
\textbf{an order of magnitude cheaper per node}, which is precisely the price of
first order and no adaptivity. At the scale of the coupled march this is the
difference between the 3-D solve's nodal update being free and the 0-D march
costing the $574$~s of Table~\ref{tab:march1um}.

\subsection{Consequence for the coupling}

\begin{table}[htb]
\centering
\caption{The two integrators side by side.}
\begin{tabular}{lll}
\toprule
& MoDELib \code{solveImmobileClusters} & ZrMicro coupled march \\
\midrule
Scheme        & semi-implicit (IMEX), nodal & CVODE BDF, adaptive \\
Order         & 1 & up to 5 \\
Step          & fixed, $5\times10^{5}$ s (20/\dpa) & adaptive ($\sim$161/\dpa\ interior) \\
Error control & none & local estimate, \code{rtol}$=10^{-6}$ \\
Stability     & L-stable on the linear losses & stiffly stable to order 5 \\
Linear algebra& none (diagonal by construction) & dense $9\times9$ LU, exact AD Jacobian \\
Cost per node & $\sim$0.14 ms / \dpa & $\sim$0.8--1.5 ms / \dpa \\
Accuracy      & first order, unquantified & $2\times10^{-6}$ measured \\
\bottomrule
\end{tabular}
\end{table}

The operator split is itself first order in $\Delta\gamma$
(Section~\ref{sec:integrators} notwithstanding, this is the scheme's design
order). MoDELib's immobile update is therefore \emph{consistent with} the
splitting order, and CVODE's six additional digits on the slow march buy nothing
unless $\Delta\gamma$ is refined alongside them.

Where the difference does matter is the return path. Once the marched field is
written back through the CD block (Section~\ref{sec:fieldcost}), the two schemes'
different truncation errors become a systematic 0-D/3-D discrepancy --- a
\emph{fourth} candidate alongside the three physics non-correspondences listed in
\code{CLAUDE.md} (the DAD bias reconciliation, the bi-pyramid vacancy family, and
size-dependent vacancy-loop emission). Anyone chasing a residual disagreement
between the interior 0-D state and the 3-D result should rule this one out
first, because it is the cheapest to test: raise \code{nSub} in
\code{solveImmobileClusters} and see whether the disagreement moves.

% ============================================================================
\section{Provenance: how the two schemes came to be}
\label{sec:provenance}
% ============================================================================

Section~\ref{sec:integrators} compared the two integrators as they stand. This
section records where each came from, because the histories explain the
asymmetry between them and carry a maintenance consequence.

\begin{table}[htb]
\centering
\caption{Immobile-species integration across the three code states.}
\label{tab:provenance}
\begin{tabular}{lccc}
\toprule
& MoDELib2-NNL & MoDELib3 & ZrMicro coupled march \\
& (upstream) & (standalone route) & (this work) \\
\midrule
\code{iSize} (immobile dofs/node) & \textbf{0} & 8 & --- \\
Immobile integrator     & \textbf{empty stub} & semi-implicit, nodal & CVODE BDF \\
\code{ImmobileSinks.h}  & \textbf{absent} & new file & n/a \\
Order                   & --- & 1 & up to 5 \\
Step control            & --- & fixed, 20/\dpa & adaptive \\
Error control           & --- & none & \code{rtol}$=10^{-6}$ \\
Loop sink on the mobile & \textbf{none} & \code{bWF\_RI}/\code{lWF\_RI} & frozen $\mathbf C_M^\star$ \\
\bottomrule
\end{tabular}
\end{table}

\subsection{The standalone scheme: from absent, to explicit, to semi-implicit}

There is no upstream scheme to compare against. In MoDELib2-NNL the immobile
machinery does not exist: \code{iSize} is zero, \code{solveImmobileClusters()}
is an empty stub, \code{ImmobileSinks.h} is absent, and the mobile species see
only the network sink through \code{otherSinks\_SI}, the source carrying a
\code{// Missing immobile sinks} placeholder where the loop term would go. The
material files corroborate this independently: \code{Zr4.txt}, the upstream-era
standalone 3-D calibration, contains \textbf{zero} immobile keys
(\code{immobileSpecies*}, \code{tauVac}, \code{nNuc}, \code{cLL},
\code{kappaLL}, \code{dad*}) against eight such key families in
\code{Zr3d\_ghoniem.txt} --- which is exactly why \code{Zr4.txt} is incompatible
with the fork. Consequently the upstream standalone route cannot produce an
\aloop/\cloop\ loop microstructure at all; the 1~\textmu m case requires
MoDELib3.

MoDELib3 introduced the field and its integrator wholesale --- but not in its
present form. It began fully explicit, and reached the semi-implicit update
through four documented corrections, recorded in
\code{MoDELib3/ZR3D\_GHONIEM\_CHANGES.md}:

\begin{table}[htb]
\centering
\caption{How the standalone immobile scheme reached its present form. Each
stiffness failure was fixed by moving one more linear loss into the implicit
denominator.}
\label{tab:evolution}
\begin{tabular}{cp{4.1cm}p{4.3cm}p{4.3cm}}
\toprule
Issue & Symptom & Cause & Fix \\
\midrule
7 & \cloop\ loops pinned on the positivity floor while \aloop\ loops
    nucleated normally
  & $\tau_{vL}$ is 35$\times$ shorter than the substep, so explicit Euler gave
    $n(1-35)<0$, clamped every iteration
  & Integrate the linear losses implicitly; returns
    $n^\star=\dot n_{\mathrm{nuc}}\tau_{vL}$ exactly as
    $\mathrm dt/\tau\to\infty$ \\
\addlinespace
8 & Thermal emission could not be made implicit
  & It was written as an absolute rate, not proportional to the stored content
  & Recast as a first-order coefficient (divide by $c$) so it can enter the
    implicit update \\
\addlinespace
14 & \aloop\ mean size 5042 at 1 \dpa, then exactly $300.0=n_{\mathrm{nuc}}$ at
     2 \dpa\ --- content collapsing $11\times$
   & The same stiffness failure in the coalescence channel left explicit:
     $\nu_{LN}\phi_{LN}\mathrm dt_{\mathrm{sub}}\approx3.6$
   & Eq.~(97) removes $n$ and $c$ proportionally, so both are first-order
     coefficients; moved implicit. \aloop\ density went from 3--4$\times$ the
     0-D to within ${\sim}30\%$ \\
\addlinespace
15 & \cloop\ content $350\times$ below the 0-D, mean size pinned at the
     nucleation size
   & Dissolution removed content proportional to the stored content
   & Release the birth content $n_{\mathrm{nuc}}n/\tau_{vL}$ instead \\
\bottomrule
\end{tabular}
\end{table}

Issue 7 is worth singling out because its diagnosis was made the right way: the
authors confirmed it was integration and not physics by checking the analytic
balance $n^\star=\dot n_{\mathrm{nuc}}\tau_{vL}=1.79\times10^{-8}$ against the
0-D $C_{vL}+C_{avL}=1.78\times10^{-8}$, then verified the fix reproduced
$1.79161\times10^{-8}$.

\subsection{The coupled march scheme: outward from the 0-D solver}

The march scheme has the opposite history. It did not have to discover
stiffness, because it inherited a stiff integrator: ZrMicro's CVODE BDF solver
predates the coupling entirely, built for the standalone 0-D 19-equation model.
The coupling was then added in three steps:

\begin{enumerate}
\item \textbf{\code{freeze\_mobile}} --- the operator split. The four mobile
      derivatives are zeroed \emph{last}, after every reaction term has already
      read them as parameters, so the same right-hand side serves as the local
      immobile micro-model with an externally supplied $\mathbf C_M^\star$.
\item \textbf{Pointwise dispatch} --- because the split makes the slow march
      decoupled between points, each quadrature point becomes one case of the
      existing OpenMP \code{--batch\_file} mode.
\item \textbf{This work} --- accumulators moved to CVODES quadrature
      ($19\to9$), the exact AD Jacobian, and per-thread workspace reuse.
\end{enumerate}

\subsection{Two routes to the same conclusion}

Both codes ended up treating the immobile loss channels implicitly, and both
were right to, but they arrived from opposite directions. MoDELib3 found the
stiffness empirically --- twice, as loop populations collapsing onto the
positivity floor --- and fixed it term by term. The 0-D side never exposed the
problem, because BDF absorbed it silently from the start; it took the spectrum
of Table~\ref{tab:spectrum} to put a number on what was being absorbed
($46\,000\times$ the forward-Euler limit).

That difference in route leaves a difference in robustness, and it is the
practical point of this section. MoDELib3's scheme is unconditionally stable on
the losses that \emph{have been} moved into \code{lossN}/\code{lossC}. It offers
no protection for one that has not. A newly added immobile loss term will be
stable only if it is first recast as a first-order coefficient and placed in the
denominator; written as an explicit rate it will reproduce issue 7 or issue 14,
and --- with no local error estimate and no step rejection --- the only symptom
will again be a species quietly pinned on the floor. CVODE would have flagged
the same term through step rejection and error-test failures without anyone
having to notice a collapsed population.

\paragraph{Sourcing.} MoDELib2-NNL is not present in this checkout: there is no
vendored copy, and \code{MoDELib3/} carries no independent git history, having
been flattened from a submodule (the six commits present are DisloCluster's
own). The upstream statements above therefore rest on MoDELib3's own change log
--- a detailed first-party record --- corroborated by the direct evidence of the
\code{Zr4.txt} versus \code{Zr3d\_ghoniem.txt} key comparison, which was checked
against the files in this repository. They have not been verified against the
upstream repository itself.

% ============================================================================
\section{Summary and recommended next steps}
% ============================================================================

\begin{table}[htb]
\centering
\caption{Summary of outcomes against intent.}
\begin{tabular}{lll}
\toprule
Change & Target metric & Outcome \\
\midrule
AD Jacobian        & residual evaluations & $-33\%$; wall neutral \\
$9\times9$ reduction & LU flops           & $-92\%$; steps $-18\%$; wall $-4\%$ \\
                   & step-size pathology  & 2048 roundoff warnings $\to$ 0 \\
Workspace reuse    & solver setup cost    & march wall $\times1.67$ at $q=1024$ \\
Field handoff      & information exchanged & 4 scalars $\to$ 289\,380 values \\
Full 1--26 \dpa\ march & same-work solver speed & $1.95\times$ (Table~\ref{tab:march1um}) \\
\bottomrule
\end{tabular}
\end{table}

\noindent
Two findings are worth carrying forward independently of the changes
themselves. First, the accumulators were quietly setting the initial step of
every coupling substep, and removing them from the error test cures a
roundoff-limited first step that had been occurring 2048 times per march
(Section~\ref{sec:pathology}). Second, the two codes integrate the same immobile
equations to very different error levels --- first order with no error control
on the 3-D side, adaptive to $10^{-6}$ on the 0-D side --- which is a systematic
0-D/3-D discrepancy distinct from the three physics non-correspondences on
record (Section~\ref{sec:integrators}).

Ordered by expected return:

\begin{enumerate}
\item \textbf{Replace the quadrature with a vector-\code{atol} state}
      (Section~\ref{sec:memo}). Keeps the step-size cure at zero extra core
      evaluations; should convert the reduction's current $-4\%$ into a clear
      win.
\item \textbf{Replace the text case-file protocol.} Every batch line repeats
      the full material parameter set; only the 19 \code{y0} values differ. At
      $q=1024$ that is ${\sim}1.5\times10^{5}$ redundant string-to-double
      conversions per substep. A base-plus-delta or binary format would attack
      what is now, after the workspace change, the largest remaining
      non-physics cost.
\item \textbf{Deduplicate quadrature points before dispatch.} In a smooth field
      neighbouring nodes integrate nearly identical initial-value problems.
      Clustering the 13-dimensional input state to a tolerance could collapse
      $q\sim10^{5}$ to $10^{2}$--$10^{3}$ genuine integrations --- likely a
      larger factor than everything above combined, and it composes with all of
      it.
\item \textbf{Guard the standalone immobile solver against the next stiff term.}
      Section~\ref{sec:provenance} shows the same failure was found twice, by
      symptom, and there is still no mechanism that would catch a third. Any new
      loss channel added to \code{solveImmobileClusters} must be recast as a
      first-order coefficient and placed in \code{lossN}/\code{lossC}; a cheap
      standing check is to run one case at \code{nSub} and at $10\,\code{nSub}$
      and assert the answers agree, which turns a silent floor collapse into a
      visible regression.
\item \textbf{Verify the 3-D fast solve end to end} in the field backend, then
      revisit GPU offload. Until the spatial solve actually runs inside the
      march loop, $q$ is set by the node count of a stored configuration rather
      than by a live coupling, and the case for batched GPU integration cannot
      be assessed on real numbers.
\end{enumerate}

\paragraph{Files.}
\code{cpp\_utils/dual.h},
\code{cpp\_utils/rate\_equations\_core.h},
\code{cpp\_utils/rate\_equations.\{h,cpp\}},
\code{cpp\_utils/jac\_check.cpp},
\code{cpp\_utils/parameters.h},
\code{code/solver.cpp},
\code{py\_utils/cpp\_bridge.py},
\code{py\_utils/modelib\_coupling.py},
\code{py\_utils/modelib\_field.py} (new),
\code{code/coupled\_0d\_3d\_ZrMicro.ipynb}.
The build now links \code{SUNDIALS::cvodes} rather than \code{SUNDIALS::cvode},
because \code{CVodeQuadInit} and its companions exist only in CVODES; the
header \code{<cvodes/cvodes.h>} was already being included.

\end{document}
